\chapter{Discussão}
\label{ch:discussao}

\section{Significado dos resultados}
\label{sec:pergunta-discussao}

O \CabotageLens{} compara duas formas de levar a mesma carga entre a mesma origem e o mesmo destino. A primeira usa rodovia direta. A segunda combina acesso rodoviário ao porto, navegação e acesso rodoviário ao destino. Essa igualdade de origem, destino e massa é necessária para que a comparação seja coerente.

O resultado não deve ser lido apenas como um total de custo ou de emissão. Ele é a última etapa de uma cadeia. Antes dele, o sistema escolhe os portos, calcula as pernas rodoviárias, identifica as viagens ANTAQ, reconstrói a carga a bordo, associa a intensidade MRV e registra a fonte de cada distância.

Os dados históricos da ANTAQ foram organizados em etapas. Uma \emph{chamada portuária} é uma ocorrência de atracação registrada na base. Chamadas consecutivas do mesmo navio no mesmo porto são reunidas em uma única \emph{parada}. Duas paradas vizinhas na sequência da viagem formam uma perna de navegação, chamada de \emph{subtrecho}. Dessa forma, 7.103 chamadas foram organizadas em 6.797 paradas pertencentes a 1.324 viagens reconstruídas. Dessas viagens foram calculados 5.438 subtrechos.

Uma mesma viagem pode fornecer vários recortes entre dois portos. Dentro da viagem real \codecell{voyage\_9612791\_00011}, o prefixo Santos--Suape--Pecém--Manaus fornece seis recortes: Santos--Suape, Santos--Pecém, Santos--Manaus, Suape--Pecém, Suape--Manaus e Pecém--Manaus. A viagem completa ainda retorna de Manaus a Santos e pode fornecer outras ligações. Os três recortes entre portos consecutivos são diretos; os demais incluem uma ou duas paradas intermediárias. Essa multiplicação explica por que 1.324 viagens geraram 12.025 recortes históricos entre um porto inicial e um porto final. No arquivo processado, cada registro recebe o nome técnico de \emph{recorte histórico viagem--OD}, em que OD significa origem--destino.

Os 12.025 recortes estão distribuídos por 177 ligações com sentido definido entre dois portos. Santos--Manaus, por exemplo, é diferente de Manaus--Santos.

Esses números mostram a quantidade de dados processados. Eles não significam que existam 177 serviços comerciais garantidos. A ANTAQ informa operações registradas no período analisado. Frequência futura, espaço disponível e contratação precisam ser confirmados fora do modelo.

\section{Quais viagens formam o indicador e qual rota fornece a distância}
\label{sec:corredores-portos-discussao}

Dois conjuntos de informação precisam permanecer separados. O primeiro é histórico e serve para formar o indicador representativo da ligação. No cálculo Santos--Manaus, o sistema aproveita a parte de cada viagem que começa na escala de Santos e termina na escala de Manaus. Entram tanto o recorte direto de \codecell{voyage\_9612789\_00004} quanto Santos--Suape--Pecém--Manaus de \codecell{voyage\_9612791\_00011}; uma passagem em que Manaus apareça antes de Santos não entra. Essa parte aproveitada da viagem pode ser direta ou incluir paradas intermediárias e recebe o nome de recorte.

Quatro verificações são aplicadas a cada recorte histórico. Primeiro, a parte aproveitada da viagem deve seguir a direção pedida. No cálculo Santos--Manaus, ela começa na escala de Santos e termina na escala de Manaus; uma sequência Manaus--Suape--Santos não serve. Segundo, nenhum dos trechos navegados entre esses portos pode estar ausente. Terceiro, deve existir uma intensidade obtida pelo IMO ou estimada e identificada como fallback. Quarto, o trabalho de transporte deve ser maior que zero para o recorte receber peso na mediana; esse trabalho é a soma da carga a bordo multiplicada pela distância de cada trecho. Um recorte com resultado igual a zero permanece registrado, mas não altera a mediana quando existem resultados positivos. O conjunto aprovado por essas regras recebe o nome de \emph{amostra de intensidade}.

O segundo conjunto pertence à execução solicitada pelo usuário. Ele contém uma única massa de carga, uma sequência de portos selecionada e a distância desse caminho. A amostra histórica ajuda a calcular o indicador típico de consumo. O caminho selecionado informa a distância sobre a qual esse indicador será aplicado à carga informada pelo usuário.

Em Santos--Manaus, a amostra possui 89 recortes históricos. Esses recortes seguem 22 sequências completas de portos. Cada sequência recebe o nome de \emph{corredor}. A amostra inclui o recorte direto. Também inclui Santos--Suape--Pecém--Manaus em \codecell{voyage\_9612791\_00011}; Santos--Rio de Janeiro--Suape--Pecém--Manaus em \codecell{voyage\_9784661\_00001}; e Santos--Navegantes--Pecém--Manaus em \codecell{voyage\_9852365\_00011}. O último não passa por Suape. Todos esses recortes ajudam a calcular um único indicador de consumo para Santos--Manaus. Esse indicador recebe o nome de \emph{intensidade representativa}.

Para escolher a distância aplicada à carga do usuário, a regra encontra um recorte direto observado. Por isso, seleciona Santos--Manaus. A matriz marítima atribui 6.112~km ou 3.300,216~nm ao subtrecho. Os portos intermediários dos outros corredores não são acrescentados a esse caminho. Eles permanecem somente na amostra estatística histórica.

Essa separação evita dois erros. O primeiro seria assumir que todo navio percorreu apenas Santos--Manaus sem escalas. O segundo seria montar uma rota artificial com partes reais de viagens diferentes. Por exemplo, Santos--Suape aparece em \codecell{voyage\_9612791\_00011}, enquanto Suape--Manaus aparece em \codecell{voyage\_9612789\_00004}; o sistema não une esses dois registros para inventar uma única viagem Santos--Suape--Manaus. Ele usa muitos recortes históricos para estimar a intensidade, mas usa uma única cadeia observada e completa para executar o cálculo da remessa do usuário.

\section{Como o sistema escolhe e registra o indicador de consumo}
\label{sec:robustez-maritima-discussao}

O indicador marítimo informa quantos gramas de combustível são consumidos para transportar uma tonelada por uma milha náutica. A unidade é g/(t$\cdot$nm), e o método chama esse indicador de \emph{intensidade}. O sistema procura primeiro o número IMO do navio no MRV. Quando encontra uma correspondência, usa o valor positivo do período mais recente para aquele navio.

Quando não encontra o IMO, o sistema estima o consumo com embarcações semelhantes. Primeiro, procura o grupo mais específico disponível, chamado de classe. Se esse grupo não fornecer um valor, usa uma categoria mais ampla, chamada de tipo, como ``navio porta-contêineres''. Essa estimativa de grupo recebe o nome de \emph{fallback} e fica identificada no resultado.

O fallback não deve ser interpretado como medição da embarcação. Para o tipo de navio, o sistema ordena os valores e retira uma quantidade inteira correspondente a 1\% de cada extremidade quando a amostra permite; caso contrário, usa a mediana. Para a classe, usa a média registrada no arquivo de eficiência depois da retirada dos valores abaixo do percentil 1 e acima do percentil 99; se ela não existir, usa a mediana. Um valor associado diretamente a um IMO não é retirado nem alterado.

No conjunto completo, 60,50\% dos subtrechos usam IMO e 39,50\% usam fallback. Em Santos--Manaus, 40 dos 89 recortes históricos usam IMO e 49 usam fallback.

Para calcular o peso estatístico de cada recorte histórico, o sistema multiplica a carga a bordo pela distância de cada subtrecho e soma os resultados. Esse total recebe o nome de \emph{trabalho de transporte}. Em Santos--Manaus, 59,54\% do trabalho de transporte está ligado a fallback. Portanto, a origem do valor precisa acompanhar o resultado. O registro dessa origem recebe o nome de \emph{proveniência}.

Para obter um único valor, o sistema ordena as intensidades da menor para a maior. Depois, acumula o trabalho de transporte nessa ordem. Escolhe a primeira intensidade em que o total acumulado alcança 50\%. Essa regra recebe o nome de \emph{mediana ponderada pelo trabalho de transporte}.

Em Santos--Manaus, a mediana ponderada é 9,322050~g/(t$\cdot$nm). A média aritmética ponderada diagnóstica é 25,243618~g/(t$\cdot$nm). A diferença ocorre porque intensidades altas deslocam a média. A mediana é menos sensível a essa parte alta da distribuição.

Isso não torna o valor certo em qualquer condição. O EU MRV resume operações reportadas em contexto europeu. Velocidade, calado, manutenção, clima, corrente, combustível e perfil de carga podem ser diferentes no Brasil. A regra robusta melhora a estabilidade do estimador, mas não elimina essas diferenças.

\section{Como interpretar a execução de uma carga}
\label{sec:porta-a-porta-discussao}

O teste de execução com 14~t apresentado no Capítulo~\ref{ch:resultados} ajuda a mostrar a função de cada dado. Na etapa histórica, as cargas reconstruídas a bordo dos 89 recortes determinam os pesos usados para obter a intensidade de 9,322050~g/(t$\cdot$nm). Essas cargas históricas não são somadas nem confundidas com a entrada controlada de 14~t.

Na execução controlada, as 14~t representam uma única entrada do modelo, e não uma carga observada na ANTAQ. A distância de 3.300,216~nm vem da matriz marítima para o caminho direto da viagem \codecell{voyage\_9612789\_00004}. A multiplicação da intensidade representativa por essa massa e por essa distância resulta em 430,71~kg de combustível de navegação.

Os 430,71~kg não representam a cadeia inteira. Para uma comparação porta a porta, ainda são necessários o percurso rodoviário até Santos, as operações portuárias representadas e o percurso rodoviário a partir de Manaus. Comparar apenas o navio com o caminhão retiraria componentes reais da alternativa multimodal.

Quando a intensidade MRV por trabalho de transporte é utilizada, o consumo anual reportado do navio é tratado como abrangente para essa fronteira. Por isso, o \emph{hoteling} não é somado novamente. Essa regra evita dupla contagem. As operações portuárias que são modeladas separadamente permanecem visíveis e devem possuir fonte identificada.

A literatura de \emph{short sea shipping} também mostra que o resultado depende do corredor, dos acessos, da utilização, da frequência, da confiabilidade e do custo total \cite{shortsea2019,modalshiftreview2020}. Essa literatura explica por que a comparação deve ser porta a porta. Ela não valida numericamente Santos--Manaus nem outro corredor brasileiro específico.

\section{Comparação com uma referência externa}
\label{sec:benchmark-discussao}

O resultado também é comparado com uma planilha produzida fora do projeto. Essa referência externa recebe o nome de \emph{benchmark}. O benchmark do Prof. Dr. Gustavo Costa permite verificar se as duas análises apontam o mesmo modo como menos emissor. Nos 21 casos com valores positivos e utilizáveis, tanto a planilha quanto o \CabotageLens{} indicam menor emissão para a alternativa com cabotagem.

Esse acordo não significa que os valores sejam iguais. A classificação \codecell{same\_direction\_large\_gap} informa que existe concordância de direção e diferença relevante de magnitude. Distâncias, carga, veículo, portos, alocação e fronteira ambiental podem ser diferentes. Antes de alinhar esses elementos, o benchmark não pode ser usado para calibrar a intensidade marítima.

\section{O que os resultados incluem e deixam de fora}
\label{sec:fronteiras-discussao}

Todo cálculo precisa declarar o que inclui e o que deixa de fora. Esse limite recebe o nome de \emph{fronteira}. Na fronteira econômica, os valores monetários incluem somente os custos operacionais representados pelo sistema. Eles ajudam a comparar os componentes modelados. Não são cotações de frete, tarifas contratadas ou prova de viabilidade comercial. Uma decisão real também depende de margem, seguro, prazo, frequência, confiabilidade, inventário, capacidade e condições do serviço \cite{competitiveness2024}.

Na fronteira ambiental, as emissões incluem somente a operação dos veículos e do navio. O relatório chama esse limite de \ttw{} de \cooe{}. Nas pernas rodoviárias, o termo específico é \emph{tank-to-wheel}. Na perna marítima, é \emph{tank-to-wake}. Um estudo \wtw{} também inclui a produção e a distribuição dos combustíveis. Uma LCA inclui outras etapas do ciclo de vida. Portanto, números \wtw{}, LCA ou de CO$_2$ isolado não podem ser comparados diretamente com a saída do sistema sem reconciliação \cite{decarb2024,maritimelca2024}.

As operações portuárias são aproximações operacionais, não inventários completos de cada terminal. A regra de não adicionar \emph{hoteling} à intensidade MRV reduz dupla contagem, mas não descreve produtividade de berço, energia auxiliar, permanência no terminal ou qualidade do ar local \cite{berth2009,shipops2022,berthairquality2010}.

\section{Valor prático e contribuição metodológica}
\label{sec:valor-pratico-triagem}

O \CabotageLens{} ajuda a escolher quais casos merecem uma análise mais profunda. Esse uso inicial recebe o nome de \emph{triagem técnica}. A ferramenta mostra quais pernas dominam o resultado, quais dados vieram de IMO, quanto do peso depende de fallback e quais premissas precisam de evidência adicional.

Um resultado favorável indica um corredor que merece investigação. O passo seguinte pode incluir consulta a armadores, terminais e transportadores, confirmação de frequência e capacidade e obtenção de tarifas. O modelo não substitui essas verificações.

Sua principal contribuição é manter cálculo e proveniência juntos. Para Santos--Manaus, por exemplo, o leitor consegue identificar os 89 recortes históricos, os 22 corredores, a divisão 40/49 entre IMO e fallback, a participação de 59,54\% do fallback, a intensidade de 9,322050 e o caminho direto de 6.112~km. Essa trilha permite explicar o resultado e também seus limites.
